\pg{350} que tem os elementos da sua entidade unidos e indivisiveis, é mais perfeito que o ente que os tem unidos mas divisiveis; porque, ao passo que o primeiro não pode perder a sua unidade e, com a unidade, a entidade, o segundo pode perder uma e outra coisa (o que é imperfeição) Logo o ente \emph{simples}, quando subsiste perfeitamente, é mais perfeito que o \emph{composto}.

Dizemos --- quando \emph{subsiste perfeitamente}; porque, se o ente \emph{simples}, ou não \emph{subsiste}, ou não subsiste \emph{perfeitamente}, mas tem razão de \emph{parte}, é menos perfeito que o \emph{composto}. Assim a alma do animal, embora simples, é comtudo menos perfeita que o animal inteiro; e a nossa alma, embora seja simples e subsistente, comtudo, por não subsistir \emph{perfeitamente} quanto á \emph{especie}, é menos perfeita que o composto humano \footnote{A \emph{simplicidade}, considerada em si, emquanto é negação de composição, não denota necessariamente nem perfeição nem imperfeição: porque, diz S. Thomaz, o que constitue a excellência do ente não é propriamente a simplicidade, mas é a perfeição, é a somma dos graus de entidade, real ou logicamente distinctos, de que o ente é constituido, Para se averiguar, pois, em concreto, se a \emph{simplicidade} importa perfeição ou imperfeição, é necessario considerar a natureza do ente, e a razão e o modo da sua simplicidade. Ora \emph{o} ente pode dizer-se e.ser simples, --- ou emquanto importa um \emph{todo completo} e \emph{subsistente} na mesma simplicidade, --- ou. emquanto importa apenas ura \emph{parte}, que não se divide noulras partes. No primeiro caso, a \emph{simplicidade} denol \emph{perfeição}, e'0 ente simples, que é um todo completo e subsistente, é inais: perfeito que o composto; porque, como dizemos no texto, conserva melhor a sua unidade e por isso a sua entidade. No segundo caso, a \emph{simplicidade} denota \emph{imperfeição} (porque quem diz \emph{parte} diz limitação) e o ente \emph{simples} é menos perfeito que o \emph{composto}.. 0 Angelico Doutor escreve: « \emph{Simplicitas} per se non est causa nobilitatis, sed \emph{perfectio}; unde ubi perfecta bonitas \emph{in} uno simplici invenitur, simplex est nobilius quam compositum: quando autem e contrario simplex est imperfectum, iunc compositum est nobilius quam simplex » (In IV \emph{Sent}., dist. 11, q. 2, a. 1 ad 1).\par Portanto a \emph{simplicidade} do ente subsistente n\&o exprime apenas uma negação de composição, mas denota uma perfeição positiva: pois denota a entidade emquanto subsiste na simplicidade. Assim como a unidade não indica apenas a negação de divisão, mas significa o proprio ente, do qual exclue a divisão; assim tambem a simplicidade significa por si a propria perfeição positiva da entidade, da qual exclue a composição. Por isso diz S. Thomaz que, no mesmo genero, um ente é tanto mais excellente, quanto é mais simples: «In quolibet genere, tanto aliquid est \emph{nobilóws}, quanto \emph{simplicius}» (\emph{C}. \emph{Gent}., \emph{l}, \emph{IS}). Ora á excellencia indica uma perfeição positiva.}.

74. O ente, que é acto puro, é absolutamente simples; o ente, que é absolutamente simples, é \textbf{acto} puro.

a) \emph{O ente}, que é \emph{acto} puro, é \emph{absolutamente simples}. --- O

1 